%%%%%%%%%%%%%%%%%%%%%%%%%%%%%%%%%%%%%%%%%
% a0poster Landscape Poster
% LaTeX Template
% Version 1.0 (22/06/13)
%
% The a0poster class was created by:
% Gerlinde Kettl and Matthias Weiser (tex@kettl.de)
%
% This template has been downloaded from:
% http://www.LaTeXTemplates.com
%
% License:
% CC BY-NC-SA 3.0 (http://creativecommons.org/licenses/by-nc-sa/3.0/)
%
%%%%%%%%%%%%%%%%%%%%%%%%%%%%%%%%%%%%%%%%%

%----------------------------------------------------------------------------------------
%	PACKAGES AND OTHER DOCUMENT CONFIGURATIONS
%----------------------------------------------------------------------------------------

\documentclass[a0, landscape]{a0poster}

\usepackage{multicol} % This is so we can have multiple columns of text side-by-side
\columnsep=100pt % This is the amount of white space between the columns in the poster
\columnseprule=3pt % This is the thickness of the black line between the columns in the poster

\usepackage[svgnames]{xcolor} % Specify colors by their 'svgnames', for a full list of all colors available see here: http://www.latextemplates.com/svgnames-colors

%\usepackage{times} % Use the times font
\usepackage{palatino} % Uncomment to use the Palatino font
\usepackage{graphicx} % Required for including images
\graphicspath{{figures/}} % Location of the graphics files
\usepackage{booktabs} % Top and bottom rules for table
\usepackage[font=small,labelfont=bf]{caption} % Required for specifying captions to tables and figures
\usepackage{amsfonts, amsmath, amsthm, amssymb} % For math fonts, symbols and environments
\usepackage{wrapfig} % Allows wrapping text around tables and figures

\begin{document}

%----------------------------------------------------------------------------------------
%	POSTER HEADER
%----------------------------------------------------------------------------------------

% The header is divided into three boxes:
% The first is 55% wide and houses the title, subtitle, names and university/organization
% The second is 25% wide and houses contact information
% The third is 19% wide and houses a logo for your university/organization or a photo of you
% The widths of these boxes can be easily edited to accommodate your content as you see fit

\begin{minipage}[b]{0.85\linewidth}
\veryHuge \color{NavyBlue} \textbf{Diffusion Kurtosis Imaging for the Human Connectome Project} \color{Black}\\ % Title
%\Huge\textit{An Exploration of Complexity}\\[1cm] % Subtitle
\huge \textbf{Ariel Rokem\textsuperscript{1}, Libby Huber\textsuperscript{2}, Parmita Mehta\textsuperscript{3}, Rafael Neto Henriques\textsuperscript{4}, Magda Balazinska\textsuperscript{3} \& Jason Yeatman\textsuperscript{2}}\\ % Author(s)
\Large 1. The eScience Institute, 2. ILABS, 3. Computer Science and Engineering, Univ. of Washington \& 4. MRC Cognition and Brain Sci., Univ. of Cambridge\\ % University/organization
\Large Contact: \texttt{arokem@uw.edu} $|$ Download: \texttt{http://arokem.org/presentations/sfn2016.html}
\end{minipage}
%
%\begin{minipage}[b]{0.25\linewidth}
%\color{DarkSlateGray}\Large \textbf{Contact Information:}\\
%Department Name\\ % Address
%University Name\\
%123 Broadway, State, Country\\\\
%Phone: +1 (000) 111 1111\\ % Phone number
%Email: \texttt{john@LaTeXTemplates.com}\\ % Email address
%\end{minipage}
%
\begin{minipage}[b]{0.19\linewidth}
\includegraphics[width=10cm]{UWlogo.png}
\includegraphics[width=8cm]{Camlogo.png}
\end{minipage}

\vspace{0.5cm} % A bit of extra whitespace between the header and poster content

%----------------------------------------------------------------------------------------

\begin{multicols}{3} % This is how many columns your poster will be broken into, a poster with many figures may benefit from less columns whereas a text-heavy poster benefits from more

%----------------------------------------------------------------------------%	Introduction
%----------------------------------------------------------------------------

\section*{Introduction}

\subsection*{The Human Connectome Project}

Diffusion MRI (dMRI) measurements provide detailed information about human brain
connectivity and microstructure \emph{in vivo}. HCP is measuring high-quality
multi-modal MRI data from 1,200 individuals and make these data publcly
available \cite{VanEssen2012}. It has already made measurements from
approximately 900 subjects available through Amazon's Simple Storage Service
(S3).

\subsection*{What model should we use to explain the data?}

\subsection*{Diffusion Tensor Imaging (DTI)}

Approximates diffusion in every voxel as a Gaussian distribution
\cite{Basser1994-hg}:

\vspace{2mm}
\begin{center}
\begin{large}

$S(\theta, b) = S_0 e^{\theta^T \mathbf{Q} \theta} $

\end{large}
\end{center}
\vspace{2mm}

\noindent $b$ is the \textbf {b-value},

\noindent $S_0$ is the signal in the absence of diffusion gradient sensitization ($b=0$)

\noindent $\mathbf{Q}$ is a positive-definite quadratic form:

\vspace{2mm}
\begin{center}

$\mathbf{D} = \begin{pmatrix} \sigma_{xx} & \sigma_{xy} & \sigma_{xz} \\
                              \sigma_{yx} & \sigma_{yy} & \sigma_{yz} \\
				                      \sigma_{zx} & \sigma_{zy} & \sigma_{zz} \\
\end{pmatrix} $

\end{center}
\vspace{2mm}

\subsubsection*{DTI: 6 parameters}

\subsection*{Diffusion Kurtosis imaging (DKI)}

DKI is an extension of DTI that accounts for non-Gaussian behavior in complex tissue, with many barriers to the diffusion process (cell membranes,
myelin sheaths, etc.) \cite{Jensen2005-vr}:

\vspace{2mm}
\begin{center}
\begin{large}

$ S(\theta, b)=S_{0}e^{-bD(\theta)+\frac{1}{6}b^{2}D(\theta)^{2}K(\theta)}$

\vspace{2mm}
\end{large}
\end{center}

\vspace{2mm}
\begin{center}
\begin{large}

$D(\theta)=\sum_{i=1}^{3}\sum_{j=1}^{3}\theta_{i}\theta_{j}Q_{ij}$

\vspace{2mm}
\end{large}
\end{center}

\vspace{2mm}
\begin{center}
\begin{large}

$K(\theta)=\frac{MD^{2}}{D(\theta)^{2}}\sum_{i=1}^{3}\sum_{j=1}^{3}\sum_{k=1}^{3} \sum_{l=1}^{3}\theta_{i}\theta_{j}\theta_{k}\theta_{l}W_{ijkl}$

\vspace{2mm}
\end{large}
\end{center}

$\mathbf{W}$ is a rank 4 tensor (3-by-3-by-3-by-3 matrix).

\subsubsection*{DKI: 21 parameters}

\subsubsection*{Diffusion Statistics}

\begin{itemize}

\item Mean diffusivity (\textbf{MD}) characterizes the mean displacement of
water molecules within a voxel.

\item Fractional anisotropy (\textbf{FA}) characterizes the variance in
diffusivity in different directions.

\end{itemize}

\noindent Variance in these parameters accounts for variance in psychiatric and
neurological disorders, normal development and aging, and psychological traits.

\large
\textbf{Brain measurements that provide a close tie between brain tissue properties and behavior}.

\subsubsection*{K-fold cross-validation}

\color{Navy}

%----------------------------------------------------------------------------%	MATERIALS AND METHODS
%----------------------------------------------------------------------------
\normalsize
\section*{Materials and Methods}

\subsection*{Data}

Measurements were obtained from the WU-Minn Human Connectome Project consortium
(\texttt{https://www.humanconnectome.org/}) through AWS S3

Of the 900 subjects that are available, 788 subjects have full diffusion
measurements: 270 diffusion weighted directions in 3 different b-values (90 directions each): 1,000, 2,000 and 3,000 $s/mm^2$

Measurements were obtained at a 1.25 x 1.25 x 1.25 $mm^3$ resolution.

A T1-weighted measurement was also obtained, and \emph{Freesurfer} was used to
segment the image into different tissue types. We focused our analysis only on
the parts of the volume that contained the white matter.

\color{DarkSlateGray}

\begin{minipage}[b]{0.75\linewidth}
  \subsection*{Dipy (Diffusion Imaging in Python)}
\end{minipage}
\begin{minipage}[b]{0.25\linewidth}
  \includegraphics[height=2cm]{dipy-logo.png}
\end{minipage}

(\texttt{http://dipy.org/})

\small

\begin{itemize}
\item Open-source software for the analysis

\item Implemented in Python. Performance is enhanced using Cython and OpenMP

\item Contains implementations of many dMRI models, including DTI, DKI and others

\end{itemize}


\begin{minipage}[b]{0.75\linewidth}
  \subsection*{Apache Spark}
\end{minipage}
\begin{minipage}[b]{0.25\linewidth}
  \includegraphics[height=3.5cm]{spark-logo.png}
\end{minipage}

(\texttt{http://spark.apache.org/})

\small

\begin{itemize}
\item Open-source parallel processing framework that enables users to run large-scale data analytics applications across clustered computers.

\item Dataflow-based execution system that provides a functional, collection-oriented API.

\item Parallel tasks are described as a directed acyclic graph (DAG) -- resilience against transient failures by tracking computational lineage

\item Optimizes for data locality when scheduling work.

\item Programming interfaces in Scala, Java, and Python.

\end{itemize}

\subsection*{Computational experiments}
We used AWS r3.2xlarge instance type, optimized for memory-intensive applications (lower price per GiB of RAM): each instance has 8 vCPU, 61 GiB Memory and 160 GB SSD storage.

\textbf{Spark}: Each compute node ran eight workers; spark default is one worker
per CPU, for a total of thirty-two workers. We used.

%------------------------------------------------

%----------------------------------------------------------------------------%	RESULTS
%----------------------------------------------------------------------------
\section*{Results}

\begin{minipage}[b]{0.5\linewidth}
  \includegraphics[width=17cm]{dti_dki.png}
\end{minipage}
\begin{minipage}[b]{0.5\linewidth}
  \includegraphics[width=17cm]{dti_1000_dki.png}
\end{minipage}

\color{SaddleBrown} % SaddleBrown color for the conclusions to make them stand out

\section*{Conclusions}
\large
\begin{itemize}

\item DKI more accurately fits the HCP data than DTI

\item DKI statistics are more reliable than DTI statistics

\item DKI has the additional benefit that it provides additional parameters that
lend themselves to a richer biophysical interpretation of the signal.

\end{itemize}

\color{DarkSlateGray} % Set the color back to DarkSlateGray for the rest of the content

%----------------------------------------------------------------------------%	Future directions
%----------------------------------------------------------------------------

\section*{Future directions}
\begin{itemize}

\item Foo

\item Bar

\end{itemize}
%---------------------------------------------------------------------------	REFERENCES
%----------------------------------------------------------------------------

\nocite{*} % Print all references regardless of whether they were cited in the poster or not
\bibliographystyle{plain} % Plain referencing style
\footnotesize \bibliography{poster} % Use the example bibliography file sample.bib

%----------------------------------------------------------------------------%	ACKNOWLEDGEMENTS
%----------------------------------------------------------------------------
\subsection*{Acknowledgements} \footnotesize This research was supported through
a grant from the Gordon \& Betty Moore Foundation and the Alfred P. Sloan
Foundation to the University of Washington eScience Institute, through NSF grant
IIS-1247469, and through gifts from Amazon and the Intel Science and Technology
Center for Big Data. We are also grateful for research cloud credits from Amazon
Web Services, and for a Google Summer of Code 2015 fellowship to RNH.

\includegraphics[height=2.5cm]{eSciencelogo.png}
\includegraphics[height=2.5cm]{SloanLogo.png}
\includegraphics[height=2.5cm]{MooreFdn.png}
\includegraphics[height=2.5cm]{NSFLogo}
%----------------------------------------------------------------------------

\end{multicols}
\end{document}
